\begin{abstract}
The deployment of Large Language Models (LLMs) as agentic orchestrators has revolutionized task automation, but the need for privacy-preserving, cost-effective solutions demands on-device inference capabilities. However, local LLMs consistently underperform compared to frontier models in tool calling scenarios, struggling with both tool selection from large tool sets and accurate argument generation for complex parameter structures. 
We introduce a methodology that disaggregates a tool-calling task into two distinct subtasks: tool selection and argument generation. We propose \emph{\decoupled}, a novel post-training approach that employs LoRA fine-tuning to create dedicated LoRA adapters for tool selection and tool-specific argument generation using separate loss masking for each of the subtasks. Furthermore, we present \sys{}, an inference framework that leverages the LoRA adapters created using \decoupled{} to perform efficient agent orchestration with the help of local models on end-user devices. \sys{} decomposes the tool-call generation step into tool selection and argument generation, and dynamically loads the corresponding LoRA adapters to generate tool calls. Additionally, \sys{} implements hierarchical orchestration to restrict the number of tools required for tool selection. Our experiments on the \mcpbench{} benchmark demonstrate that the \baselinemodel{} model trained using \decoupled{} improves the tool calling accuracy of the base model by 46\%, and outperforms other local reasoning, non-reasoning and fine-tuned models of similar size in all cases, and models that are 2\myx larger, in most cases.
\end{abstract}
%-------------------------------------------------------------------------------
\section{Introduction}
%-------------------------------------------------------------------------------

The emergence of Large Language Models (LLMs) has enabled the development of sophisticated agent systems that can interpret natural language commands and coordinate multiple tools to execute complex, end-to-end tasks. These agent systems use an LLM as the central controller (also called as orchestrator) that receives user requests and a list of available tools. The orchestrator then works through the task step-by-step: it calls a tool, stores the results, uses those results to decide the next action, calls another tool, and builds up a collection of information from each step. The process continues until the controller has gathered enough data to provide a complete response to the user's original request.

The proliferation of agentic architectures has driven the establishment of standardized integration protocols, most notably the Model Context Protocol (MCP) \cite{anthropic_mcp_2024}. MCP establishes a standardized protocol for LLMs to communicate with external applications. Every MCP application defines a set of tools (called as a ``toolset'' or MCP Server), that contains all the tools belonging to the application, along with their descriptions and instructions on how the LLM can use them. Leveraging this protocol, practitioners have built a vibrant ecosystem, from reading and writing local files to invoking remote APIs \cite{mcp_servers_2025}, all of which are accessible directly through LLMs. 
% \keisuke{give definition of toolset, MCP server here}

Nevertheless, incorporating applications into agentic frameworks raises security concerns, as it can involve exposing sensitive or private information residing within these applications to LLMs \cite{hasan2025model, hou2025model, guo2025systematic}. This makes on-device inference critical for agentic systems. Running LLMs locally on end-user devices enhances privacy by ensuring that personal data remains on-device, while simultaneously eliminating the costly API expenses associated with orchestration via frontier models \cite{gu2025consumerbench}. However, this shift toward local deployments raises a critical question: Can locally deployed LLMs match the orchestration performance of their cloud-based counterparts?

Our investigation using MCP evaluation benchmarks \cite{wang2025mcp} reveals two core shortcomings that make existing local LLMs ineffective as agentic orchestrators. First, these models demonstrate \textbf{poor tool selection capabilities}, frequently choosing inappropriate tools due to limited ability to reason about the right tool usage from sometimes ambiguous descriptions. This challenge intensifies with larger tool sets, as expanded context length (often exceeding tens of thousands of tokens) overwhelms local models' attention mechanisms. Second, local LLMs exhibit \textbf{poor argument generation capabilities}, frequently failing to produce accurate parameters for tools with structural requirements. They also have limited ability to fix the mistakes in subsequent steps, causing repeated failures in tool calling. Furthermore, our experiments show that straightforward approaches such as prompt tuning yield marginal accuracy improvements, but fail to improve local models' orchestration capabilities. 

A popular approach for improving the tool-calling capability of local models is fine-tuning~\cite{erdogan2024tinyagent,lin2024hammer,liu2024toolace}. While traditional fine-tuning helps in improving the accuracy of local models for specialized tasks, its improvements in tool-calling performance are not substantial. This is because traditional fine-tuning requires that the local models simultaneously learn both tool selection and argument generation. Fundamentally, these two capabilities differ from each other in that tool selection is a classification task (identifying appropriate tools from available options), while argument generation combines task-specific parameter selection with syntactic accuracy.

To overcome these constraints, we introduce a novel methodology that disaggregates the tool-calling task into two distinct specialized components: (1) \textbf{tool selection} and (2) \textbf{argument generation}. Leveraging this decomposition, we propose \textbf{\decoupled}, a post-training approach that employs LoRA fine-tuning \cite{hu2022lora} to create dedicated adapters for tool selection and argument generation for individual tools, separately. \Decoupled{} uses separate loss masking for tool selection and argument generation, and creates a dedicated LoRA adapter for each tool to generate its arguments, and a common LoRA adapter for the classification task of selecting the right tool at every step. We build an automated pipeline including synthetic data generation to produce diverse and high-quality training datasets and performing \decoupled{} on the dataset to create the LoRA adapters.

This paper presents \sys{}, an inference framework for local model orchestration that leverages the LoRA adapters generated through \decoupled{} to perform efficient agent orchestration with the help of local models on end-user devices. \sys{} breaks down every generation step into tool selection and argument generation, and dynamically loads the corresponding adapters for generating a tool call. Furthermore, \sys{} implements \emph{hierarchical orchestration} that enables it to support a large number of tools without compromising on accuracy. Hierarchical orchestration employs a second layer of decoupling that uses the base model to perform the high-level routing task of selecting the most appropriate toolset (such as \texttt{filesystem} or \texttt{notion}) at every step, which then routes the request to the tool selector that is trained on the particular toolset. This helps the tool selectors to choose from a restricted set of tools and thereby use a smaller context length containing only the restricted set of tools. 

We use \decoupled{} on the baseline \baselinemodel{} to create an LLM orchestrator called \sysmodel{}. Extensive experimentation across two benchmarks validates \sys{}'s effectiveness. Results demonstrate that \sysmodel{} significantly enhances tool orchestration quality compared to other local off-the-shelf and fine-tuned models of a similar size. \sys{} also performs similar or better than local reasoning models, thus achieving high accuracy along with lower latency. Additionally, \decoupled{} delivers 2\myx higher improvement in tool calling compared to baseline on a standard MCP benchmark for the \texttt{filesystem} toolset, compared to traditional fine-tuning on the same dataset. \sys{} operates efficiently on consumer-grade hardware, democratizing access to powerful and privacy-preserving agentic AI systems.

\section{Analysis}

In this section, we analyze the effectiveness of local LLMs as orchestrators for agentic systems. To this end, we developed an agentic system in Rust that operates in a loop, beginning with a system prompt. It calls an LLM orchestrator to generate a tool call at each step. It then executes this tool in a containerized environment 
% \sean{maybe we can briefly discuss the choise of why tool call execution is happening in a containerized environment (e.g., each tool call system access can be monitored)} 
and appends the outcome to the conversational context, continuing to the next iteration. The workflow concludes when the orchestrator determines the task is complete and generates no further tool calls.

\subsection{Evaluation Benchmark and Metric}
\label{analysis:benchmark_metric}
\input{tables/tbl-analysis-perf}

\paragraph{Benchmark}
We use MCP-Bench~\cite{wang2025mcp}, a recently proposed benchmark that evaluates LLM
agents in realistic tool-use scenarios. Tasks in MCP-Bench are generated using an LLM-based synthesis pipeline which uses tool I/O signatures to create dependency chains of tool invocations, and then translates these dependency chains into natural language instructions, called ``task descriptions''. This is followed by a query-fuzzing phase where each task is rewritten such that it retains the core objective but omits explicit tool references and execution steps. An example of the task description and the fuzzed query is in \sref{appendix:task_description}. We use MCP-Bench to create 50 tasks for the \texttt{filesystem} MCP toolset~\cite{mcp_filesystem_server}, a collection of 12 tools for managing files and directories such as reading and writing files, creating directories, and querying metadata such as file size.

\paragraph{Metric}
To ensure objective and scalable evaluation, we employ an LLM-as-judge methodology \cite{zheng2023judging}. We use GPT-5 as the judge, which is provided with the ground-truth state of the filesystem for each task, allowing for a deterministic and accurate assessment of the final outcome.
We assess orchestration performance using \emph{\judgemetric}, a metric which evaluates the output of agent systems on a 0-10 scale for tool-calling proficiency. The evaluation framework decomposes user queries and leverages the ground-truth data to establish an expected set of tool invocations that the agent system should produce. A model's score represents the proportion of expected tool outputs successfully generated relative to the expected tool calls. Note that there can be more than one tool that is suitable for a particular task; the judge is instructed to consider all such suitable tools as valid in \sref{appendix:toolfit_prompt}. 
% \sean{two questions: 1. what if a task can be done with the same tool calls but different orders, in this case do we only accept the predicted sequence completely match the ground truth sequence (e.g.,  task: create a node project, one can create all empty files first and then start to write their contents given filesystem MCP or complete each file first before going to next, both are corrects), 2. what if there are two totally different valid tool call sequence (e.g., task: given a dir summarize the files within it, one can call readfile for each file or call readmultiplefiles(path))} 
A perfect score of 10 indicates complete alignment, where the model correctly invokes all anticipated tools.

% It is important to note that this evaluation focuses exclusively on tool-calling competency rather than response synthesis\baris{unclear what this means here as you didn't really define this well and distinguish it and didn't explain how the synthesis is done.} quality. We exclude the final summarization phase from the agentic workflow and apply our evaluation framework to the complete tool execution sequence, enabling objective and systematic assessment across different models.  

% We assess each model's performance across two key dimensions:
% \begin{itemize}
%     \item \textbf{Tool Appropriateness:} Whether the model selected the correct tool at each step.
%     \item \textbf{Tool Argument Quality:} Whether the model provided correct and well-formed arguments for its chosen tools.
% \end{itemize}
% Each dimension is scored on a scale from 0 to 10, where 10 signifies a perfect outcome. The specific system prompts used for scoring are available in the Appendix. From these scores, we derive a binary \textbf{Task Success} metric, which is set to ``1'' if Task Fulfillment is 7 or higher, and ``0'' otherwise.

\subsection{Local Models Evaluated}
Based on the typical memory capacity of consumer-grade GPUs, we consider models smaller than 24GB in size for model weights as local models, so as to fit the model weights as well as the key-value cache in the GPU memory. We test a variety of these local models on the benchmark, as specified below:

\paragraph{Off-the-shelf Local Models}
 These include \basemodel{} \cite{yang2024qwen25}, \llamathreeeightb{} \cite{grattafiori2024llama}, \qwenthreeeightb{} \cite{yang2025qwen3} (with and without reasoning tokens) and Qwen-3-32B-Quant~\cite{yang2025qwen3}

\paragraph{Local Fine-tuned Models}
These include popular fine-tuned function-calling models, namely, \xlam{}~\cite{zhang2024xlam} and \toolace{}~\cite{liu2024toolace}. 

\subsection{Local Models are incapable as LLM orchestrators}
As shown in \Cref{tbl:analysis:perf}, our evaluation reveals that existing local non-reasoning models perform significantly worse than the frontier \gptfivemini{} model. The closest top-performing local model, Qwen-3-32B-Quant is a reasoning model, which incurs a high latency of generation, a known problem with the reasoning-based models~\cite{liu2025efficientinferencelargereasoning}. This shows that local models are not able to perform effectively as orchestrators for agentic systems. 

\paragraph{Prompt Tuning Is Not the Solution}
We next investigate whether a prompt-tuning-like approach can be used to significantly improve the performance of local models. While we cannot validate the effectiveness of prompt tuning exhaustively, we check the performance of \baselinemodel{} for the straightforward, non-fuzzed task descriptions (explained in \sref{analysis:benchmark_metric}) for the same 50 tasks. These tasks contain explicit instructions about tool names and tool arguments. We observe that even using explicit task descriptions incrementally improves the performance of \baselinemodel{}, from 16\% to 22\%. We show this impact of prompt tuning on other local models in \sref{appendix:impact_of_prompt_tuning}.


% The top-performing local model, \todo{model}, only achieved a \judgemetric{} score of \todo{??}. 
% A closer look at the metrics reveals two primary failure modes. First, in \todo{??} of the failed cases, the \judgemetric{} score was low, indicating the model struggled to select the correct tool. Second, even when the right tool was chosen, \todo{??}\% of these instances received a low \judgemetric{} score, showing a distinct difficulty in generating proper arguments.

\input{tables/tbl-analysis-ablation}

\subsection{Problems in Tool Selection and Argument Generation}
To pinpoint the sources of failure in local models, we conducted an ablation study that decouples the two key stages of tool use: tool selection and argument generation. We achieve this by modifying the agent's workflow to use two separate LLM calls at each step: one to generate the tool's name and a second to generate its arguments.

This separation allows us to isolate each stage by substituting a frontier model's output for one of the steps. For example, to measure the quality of a local model's tool selection, we use it to choose the tool name but rely on \gptfivemini{} to generate the arguments. Conversely, to assess its argument generation capability, we use \gptfivemini{} to select the tool and the local model to generate the arguments. The results of this breakdown for the \baselinemodel{} local model are presented in \Cref{tbl:analysis:ablation}.

% We observe that when we delegate tool selection to a powerful frontier model, \baselinemodel{} improves by 2\myx, which shows its incapability to select the right tools for a task. Similarly, delegating argument generation also results in a similar improvement, which also shows the incapability of \baselinemodel{} in generating the right arguments. 

We find that the performance bottleneck is primarily for the tool selection as changing the tool selector to \gptfivemini{} improves the \judgemetric{} from 16\% to 60.8\%. On the other hand, delegating the argument generation to \gptfivemini{} only improves the performance from 16\% to 28.8\%. More importantly, this performance breakdown also suggests that tool selection and argument generation fundamentally differ in their nature, and can benefit from different optimizations. Fundamentally, tool selection is a classification problem, which requires identifying the appropriate tool at every step in the agent workflow from the available options. Argument generation, on the other hand, requires generating the right arguments as well as producing a structured output with syntactic accuracy. 

\subsection{Problems with Long Context due to Tool Descriptions}

The attention mechanism of local models is often impacted due to long context inputs, as identified by prior work~\cite{liu-etal-2024-lost, qin-etal-2023-nlp}. Even when the total number of tokens is smaller than the maximum context length that the model supports, a large context increases the complexity of the attention mechanism, causing the accuracy to drop. Furthermore, it also increases the latency of generation -- the reasoning models (\qwenthreeeightb{}, Qwen-3-32B-Quant) require 10-20\myx higher end-to-end time compared to the non-reasoning local models on the consumer-grade RTX 6000 GPU, for processing the queries in \mcpbench. 

We run the same benchmark on \baselinemodel{}, but we increase the list of available tools to include tools from 2 other common applications: Notion and monday.com. This increases the total number of tools from 12 to 30. MCP mandates every tool to provide a description of its functionality, along with a json schema that contains the required and optional arguments for the tool. As a result of this, the tool descriptions account for a total of 9.8K tokens, as opposed to only 3K-4K tokens belonging to the actual task to be performed. We find that the increased context length reduces the \judgemetric{} for \basemodel{} from 16\% to 10.4\%.

% achieve a high \judgemetric{}, showing its incapability to generate the right tool arguments. Similarly, delegating the argument generation task to the powerful \gptfivemini{} also does not result in a significant improvement 




% To investigate if prompt engineering could mitigate these issues, we created a simplified version of the prompts that removes the fuzzing of the prompts, to create a more straightforward sequence of steps for the local models to execute. While this modification led to a marginal performance improvement, the quality gap between local and frontier models remained substantial. This finding suggests that the performance deficit is due to fundamental capability limitations in local models rather than insufficient prompt clarity. \keisuke{show example of the modified prompts possibly in Appendix}\baris{indeed, more clarity is needed here. since we don't really know what fuzzing is, we can't be convinced that turning it off simplifies this. Need more context and explanation to make this text make sense to readers.}

% \subsection{Performance Breakdown}
% To pinpoint the sources of failure in local models, we conducted an ablation study that decouples the two key stages of tool use: tool selection and argument generation. We achieve this by modifying the agent's workflow to use two separate LLM calls at each step: one to generate the tool's name and a second to generate its arguments.

% This separation allows us to isolate each stage by substituting a frontier model's output for one of the steps. For example, to measure the quality of a local model's tool selection, we use it to choose the tool name but rely on gpt-5-mini to generate the arguments. Conversely, to assess its argument generation capability, we use gpt-5-mini to select the tool and the local model to generate the arguments. The results of this breakdown for the best-performing local model are presented in \Cref{??}.

% \todo{Draw insights from this to observe why tool selection is bad and why argument generation is bad. }



% \subsection{Takeaways}

% \todo{Come up with takeaways that motivate the design: such as why tool selection is bad and why tool arguments are bad. How much does context length affect, etc. They should be backed by data.}

\onecolumn
\section*{Appendices}
\appendix

\section{Impact of Prompt Tuning}
\label{appendix:impact_of_prompt_tuning}
\input{tables/tbl-appendix-prompttuning}

\section{Prompt for Query Generation}
\label{appendix:prompt_for_query_generation}
The following passage shows the exact prompt used for the generation of file system queries: 
% (Assumes you've defined ColorBlock as before)

\begin{ColorBlock}[Prompt]
\ttfamily\small
\textbf{SYSTEM PROMPT:} You are generating ONE English user request (exactly one lines, no quotes, no code block, no numbering, no explanations).\\
\textbf{GOAL:} Create ONE requests that need to call \$TOOL\_NAME.\\
\textbf{OUTPUT RULES (must follow all):}\\
\begin{itemize}[leftmargin=*, itemsep=2pt]
  \item Output EXACTLY TEN lines containing ONLY the user’s request.
    \begin{itemize}[itemsep=1pt]
      \item No prefixes/suffixes
      \item No labels
      \item No extra lines
    \end{itemize}

  \item English only.

  \item Never include absolute paths.
    \begin{itemize}[itemsep=1pt]
      \item Refer to location generically (e.g., “in the allowed directory”, “within the permitted workspace”, “inside the sandboxed area”)
      \item Or use a relative subpath (e.g., reports/2024-Q3.csv)
      \item Do not include drive letters
      \item Do not include leading “/” root paths
    \end{itemize}

  \item Vary phrasing and structure aggressively:
    \begin{itemize}[itemsep=1pt]
      \item Imperative / interrogative
      \item Polite / terse / conditional
      \item With or without “please”
      \item Passive / active voice
      \item Different synonyms for “allowed directory”
    \end{itemize}

  \item Use diverse item names and extensions:
    \begin{itemize}[itemsep=1pt]
      \item logs, csv, json, txt, md
      \item pdf, png, jpg, mp3, mp4
      \item zip, tar.gz, .env
      \item Hidden dotfiles
      \item Names with spaces/Unicode/UPPERCASE
      \item Nested relative subfolders
    \end{itemize}

  \item Include concrete numbers:
    \begin{itemize}[itemsep=1pt]
      \item “first 10 lines”, “last 5 lines”
      \item “image1.png”, “report.pdf”
      \item “10 MB”
    \end{itemize}

  \item Do not mention any tool names, schemas, or parameters explicitly.
\end{itemize}
The global list of tools is as follows: \$GLOBAL\_TOOL\_LIST
\end{ColorBlock}

\section{An Example of Task Description and Fuzzed Query}
\label{appendix:task_description}
\begin{ColorBlock}[Task Description and Fuzzed Query]
\ttfamily\small
\textbf{Task Description: }%
\begin{enumerate}
  \item Verify allowed directories.

  \item Obtain a recursive directory tree for \texttt{rena-browserd/browserd/src}.

  \item For exactly these four subdirectories, list entries with sizes and sort by size (\texttt{sortBy=size}):
    \begin{itemize}
      \item \texttt{rena-browserd/browserd/src/inference\_engine}
      \item \texttt{rena-browserd/browserd/src/container}
      \item \texttt{rena-browserd/browserd/src/container\_comm}
      \item \texttt{rena-browserd/browserd/src/app\_registry}
    \end{itemize}

  \item From those four listings, identify the three individual files with the largest sizes that have a \texttt{.rs} extension.

  \item For each of the top-3 files:
    \begin{itemize}
      \item Read the first 50 lines (\texttt{head=50}).
      \item Check whether a header is present, defined as: at least two of the first ten lines start with \texttt{//} or begin a block comment \texttt{/*}.
    \end{itemize}

  \item Fetch file metadata for each top file:
    \begin{itemize}
      \item Size
      \item Last modified
      \item Permissions
    \end{itemize}

  \item Produce a JSON report (array), one object per file containing:
    \begin{itemize}
      \item \texttt{path}
      \item \texttt{size\_bytes}
      \item \texttt{header\_present} (true/false)
      \item \texttt{last\_modified\_iso}
      \item \texttt{needs\_deep\_audit} (true if \texttt{size\_bytes > 15000} \textbf{OR} last\_modified is within the past 7 days)
      \item Short recommendation string
    \end{itemize}
\end{enumerate}

% Constraints / notes
\begin{itemize}
  \item Explicitly traverse only the four named subdirectories (no other directories).
  \item Process exactly the three identified files (no more than three \texttt{read\_file}/\texttt{get\_file\_info} calls each).
  \item Treat relative dates as: past 7 days and past 90 days when computing recency.
\end{itemize}

\textbf{Fuzzed Query: }%
I'm trying to get a repo ready for a security review for my project and my boss only wants me to look under \path{rena-browserd/browserd/src} — specifically those four subfolders \path{rena-browserd/browserd/src/container}, \path{rena-browserd/browserd/src/app_registry}, \path{rena-browserd/browserd/src/inference_engine}, and \path{rena-browserd/browserd/src/container_comm} — and I could use help figuring out which three individual .rs files are the largest by file size: can you list the entries in each of those four subfolders with sizes sorted by size so we can pick the top three .rs files, then read the first 50 lines of each selected file to check whether a header seems present (define header as at least two of the first ten lines starting with \texttt{//} or \texttt{/*}), fetch each file's size, last modified time, and permissions, and return a JSON array (with paths given as relative paths from the root permitted directory) where each object has path, size\_bytes, header\_present (true/false), last\_modified\_iso, needs\_deep\_audit (true if size\_bytes > 15000 or last modified within the past 7 days), and a short recommendation …

\end{ColorBlock}


\section{Prompt Used for Computing ToolFit}
\label{appendix:toolfit_prompt}
\begin{ColorBlock}[Prompt]
\ttfamily\small
You are an impartial Judger that evaluates whether a tool-calling trajectory adequately satisfied a filesystem-related query using only the tools invoked and their raw outputs, not any assistant summaries.\\
\textbf{Your job is to compute coverage:}
\begin{itemize}
    \item Get the user query requirements from the query, for the output
    \item Map each user query requirement to tool requirements which can help achieve that requirement (every user query requirement must necessarily have one tool that is relevant to it). This can be fetched by looking at tool descriptions
    \item Check if the tool output exists in the user trajectory
    item Score the trajector from 0-10 based on how many tool outputs exist, as a percentage of the required tool outputs
\end{itemize}
\textbf{Inputs:}You will be given four inputs:
\begin{itemize}
    \item input\_a (fs\_status): the ground-truth snapshot of the relevant filesystem (directory listings and/or metadata).
    \item input\_b: (tool descriptions): This contains the descriptions of all the tools that were available for the trajectory. 
    \item input\_c (query): the user's natural-language request (e.g., list files, read file \path{contents/data}, fetch metadata like \path{size/mtime/ctime/permissions}, produce JSON, etc.).
    \item input\_d (tool\_call + tool\_response): the exact tools invoked and their raw outputs. This is the only evidence of what the agent retrieved. Assume there is no final assistant answer.
\end{itemize}
\textbf{Getting user query requirements:}
\begin{itemize}
  \item Define the set of atomic requirements implied by the query:
    \begin{itemize}
      \item Listing queries: one atomic requirement per relevant item that should appear (file/dir).
      \item Metadata queries: one atomic requirement per (item, requested\_field) pair (e.g., (fileA, mtime), (fileA, size)).
      \item Content/data queries: one atomic requirement per item whose contents are requested. Treat as satisfied only if contents are shown and not truncated such that the task can be completed. If explicit ranges are requested (e.g., first 10 lines), treat each requested range as its own atomic requirement.
      \item ``All''/pattern scope: expand the item set using \texttt{fs\_status} (e.g., all files under a directory or matching an extension/glob). Missing any item in scope yields missing atomic tool requirements for that item.
      \item When uncertain about presence/completeness, or when elements are only implied (not evidenced in tool outputs), treat those atomic requirements as unsatisfied.
    \end{itemize}
\end{itemize}

\textbf{Mapping user query requirements to tool requirements:}
\begin{itemize}
  \item Based on the tool description and \texttt{fs\_status}, map each user query requirement to the tool that can achieve it. There MUST BE one / more tools that can help achieve EVERY user query requirement. If there are multiple tools that can achieve the user query requirements, then having any of them is valid.
\end{itemize}
\textbf{Your Tasks}
\begin{itemize}
  \item List relevant files/directories. From the query, derive a concrete list of relevant files/directories (explicit paths, or those implied by patterns/globs). If none, use []. Prefer exact paths; include directories if applicable.

  \item Derive the tool requirements. Using the user requirements, tool descriptions and \texttt{fs\_status}, enumerate all atomic tool requirements as defined above.

  \item Evidence check (numerator). Using only \texttt{tool\_response}:
    \begin{itemize}
      \item Mark an atomic requirement satisfied if the tool requirement item appears and the needed data for that atom is present (for metadata: the specific field value; for content: the actual content, not truncated; for listings: the item name/path).
      \item If content is clearly truncated (e.g., elided, "output too long"), mark the content requirement unsatisfied unless the query asked only for a subset that is fully present.
      \item If fields are missing, wrong, or unverifiable, mark unsatisfied.
      \item Extra/irrelevant outputs do not count against coverage; simply ignore them.
    \end{itemize}

  \item Compute coverage.
    \begin{itemize}
      \item \texttt{coverage\_percent = (satisfied\_atomic\_requirements / total\_atomic\_tool\_requirements) * 100}
      \item If \texttt{total\_atomic\_tool\_requirements == 0}, set coverage to 0\% (cannot verify anything) and explain.
    \end{itemize}

  \item Map to score (0--10).
    \begin{itemize}
      \item \texttt{Score\_ToolCoverage = round(coverage\_percent / 10)} producing an integer from 0 to 10 (0\%→0, 100\%→10; 95\%→10, 94\%→9, etc.). Use standard rounding (.5 rounds up).
    \end{itemize}

  \item Output Format (strict JSON, no code fences, no extra keys)
    \begin{itemize}
      \item Produce exactly:
\begin{JSONblock}
{
"Reasoning_ToolCoverage": "<one concise paragraph that: (a) lists the relevant files/directories inline (e.g., [\"/path/a.txt\", \"/logs/b.csv\"]) or []; (b) summarizes the atomic tool requirements and how many were satisfied vs total; (c) states whether required contents/metadata/JSON values were fully evidenced or truncated/missing; (d) notes any missing scope items or fields that reduced coverage.>",
"Score_ToolCoverage": <integer 0 to 10>
}
\end{JSONblock}
      \item Reasoning\_ToolCoverage must be a single paragraph.
      \item Score\_ToolCoverage must be an integer between 0 and 10.
    \end{itemize}
\end{itemize}

\textbf{Additional Principles}
\begin{itemize}
  \item Judge only from tool outputs and \texttt{fs\_status}; never infer unseen data or rely on assistant text.
  \item For "all"/pattern queries, completeness of scope is derived from \texttt{fs\_status}.
  \item If any required element is unclear, unverifiable, or implied only by context, treat it as unsatisfied.
  \item Do not penalize formatting; coverage concerns presence and completeness of required data.
\end{itemize}
\end{ColorBlock}
% \vfill
% \section{Decoupled Fine-Tuning using \sys{}}

% This paper proposes \sys{}, a fine-tuning pipeline for enabling practical adoption of local LLMs for serving as orchestrators for agentic systems. The main design principle behind \sys{} is ``\decoupled{}'', which decouples the task of tool calling into tool selection and argument generation, and separately fine-tunes each of these tasks using LoRA, using the same fine-tuning dataset. This decoupling allows \sys{} to create a separate LoRA adapter for each tool to generate its arguments, along with a separate adapter to act as a classifier, to select a tool name (and thus, the LoRA adapter) at each step to generate a tool call. Furthermore, \sys{} uses hierarchical orchestration, which is yet another decoupling, for reducing the context of the classifier, to work with a reduced set of tool options which it has been fine-tuned on. 
% \Cref{??} shows an overview of the \sys{} pipeline. 

% In this section, we describe the fine-tuning approach used in \sys{} that includes data generation for training and function masking for computing the loss. We then describe our approach for hierarchical orchestration to reduce the context length during inference, and finally discuss the integration of all these techniques within an inference framework to achieve effective local orchestration.   

% \subsection{Fine-tuning pipeline}
% The goal in \decoupled{} is to fine-tune a local model to behave as a classifier, for selecting the appropriate tool to call at each step from the available tools, and to also use a separate adapter for every tool, to generate the right arguments for the tool-call. In order to achieve this, it is important that the training dataset contains a roughly equal proportion of prompts for every tool from the available set of tools, to avoid biasing the classifier towards a particular tool, and such that every tool-specific adapter has enough examples to learn its tool parameters. 

% \subsubsection{Dataset Generation}
% In order to automate and scale the process of prompt generation for fine-tuning, we use the frontier gpt-5-mini model to general the training data. We instruct the frontier model to generate 1,000 training examples for each tool in the available set of tools, where the prompts include a task requiring the invocation of that tool. The Appendix contains the prompt that we use for generating the fine-tuning dataset for one of the tools belonging to the filesystem MCP server.

% After generating the prompts, we use the frontier model to solve the prompts using our agentic framework system, with gpt-5-mini as the orchestrator LLM, and log the entire trajectories, which include the prompt, LLM generation steps and tool outputs. Once we get the trajectories for all the tools for an MCP server, we divide it into a training set and validation set of 80-20, and use it for fine-tuning the local model for classification. 

% In order to get data for classifying the tool adapters, we separate the trajectories based on the tools that they invoke.

% \subsubsection{Loss Masking}

% While fine-tuning the classifier model, we apply a loss masking function, that masks the arguments of a tool, and simply computes a loss over the tool name. 

% The loss masking for tool-specific adapters works differently. \todo{Write about this @Zhan}.

% \subsection{Hierarchical Orchestration}
% Once we have a fine-tuned classifier for selecting the right tool at every step from the available tools of an MCP server, we generalize our approach to support multiple MCP servers via hierarchical orchestration, to contain the context length of the model. The way this is done is to identify the appropriate MCP server to use at every step, and then to invoke the specific classifier at that step. We find that this task does not require fine-tuning, but rather can be done by the base local model. We use structured decoding and a system prompt to instruct the base model to output the MCP server to use at every step, which then uses the MCP-server-specific classifier and tool-call adapters for that step which operate on a reduced context. 

% \subsection{Putting it all together}

% Finally, we put it all together in an agent orchestration framework called \sys{}. \sys{} is written in Rust, and uses containerized environments for invoking tools, for fault isolation and resource management. We now explain how \sys{} works as an agentic system, using local orchestration. 

% We use vLLM as our inference backend, that is loaded with an adapter that acts as a classifier for each MCP server, and a separate adapter for each tool belonging to each server. Note that we use LoRA adapters, and vLLM can load and unload hundreds of LoRA adapters, making this approach practical. 

% On getting a user prompt, \sys{} operates in a loop fashion. At every step, \sys{} first performs hierarchical orchestration, where it selects the MCP server to use for that step. Then \sys{} makes a inference call to the MCP-server-specific classifier to select the right tool, and then a call to the tool-specific adapter, to generate the tool call. This tool call is then executed to get the tool output, which gets added to the context for the next step. This loop continues until the MCP-server-specific classifier outputs a special tool-call called "summarization", which instructs the local model to generate a summary of the content based on the user input and ends execution. 
\section{Efficient Local Agent Orchestration with \sys{}}

We introduce \sys{}, a novel framework designed to enable the practical adoption and deployment of local LLMs in agentic systems via Low-Rank Adaptation \cite{hu2022lora}. The core design principle of \sys{} is \textbf{\decoupled}, a strategy that decomposes the complex agent workflow at two levels. 

At the first level, it separates tool selection from argument generation. This decoupling allows the creation of a dedicated LoRA adapters for the tool selection process and for the argument generation process. The output of \decoupled{} is a single tool selection adapter and multiple argument generation adapters. The tool selection adapter is responsible for performing the classification task of choosing the right tool to invoke at every step, and every tool has a dedicated argument generation adapter that generates the tool-call arguments for its tool, given an input conversation history. 

\sys{} employs a second level of decoupling through \emph{hierarchical orchestration}, which enables local models to operate on shorter context lengths, thereby benefiting both accuracy and latency. This strategy allows using a separate tool selector for every toolset (e.g., \texttt{filesystem} or \texttt{notion}), thus limiting the tool selector's decision space and improving its accuracy.

In this section, we detail the \sys{} methodology. We begin by describing the fine-tuning process that includes our synthetic data generation strategy and the use of separate loss masking to train the tool selector and argument generator adapters. We then explain the hierarchical orchestration mechanism for managing a large number of tools. Finally, we discuss the integration of these components into a robust inference framework for effective local orchestration. 

% \subsection{The \Decoupled{} Pipeline}
\subsection{Decoupled Fine-Tuning Pipeline}
The primary objective of the \decoupled{} is to enhance the tool-calling accuracy of a general-purpose LLMs. This is achieved by fine-tuning two distinct types of LoRA 
% The primary objective of the \decoupled{} pipeline is to transform a general-purpose local LLM into a highly specialized tool-calling engine. This is achieved by fine-tuning two distinct types of LoRA
adapters. First, we train a \emph{Tool Selector} adapter that identifies the most appropriate tool among a given toolset to use at each step of a task. It effectively acts as a router, outputting only the name of the tool to be used. Second, we also train a separate \emph{Argument Generator} adapter for each individual tool in the toolset. The sole responsibility of this adapter is to generate the correct and well-formed arguments required by its corresponding tool.

A critical requirement for this approach is a high-quality and diverse training dataset. To prevent the classifier from being biased towards frequently used tools and to ensure each argument adapter has sufficient data to learn its specific parameter schema, the dataset must contain a rich and evenly distributed set of examples for every available tool.

\subsubsection{Synthetic Dataset Generation}
To automate and scale the creation of our fine-tuning dataset, we leverage a powerful frontier model, \gptfivemini{}, as a data generator. For each tool in a given toolset (\eg, \texttt{filesystem}), we prompt \gptfivemini{} to generate 1,000 diverse training examples, where each example consists of a user-like prompt that necessitates the use of that specific tool. The prompt used for this data synthesis process is detailed in the Appendix \sref{appendix:prompt_for_query_generation}.
% \sean {maybe we can provide a definition of what is a tool what is a mcp server, what is the input output json schema earlier, it helps to understand the paper easier}\baris{yes even in the intro and at the begining of design, give an example first time use use the terms so it is clear.}..}

After generating the task prompts, we use the same frontier model (\gptfivemini{}) to generate the tool-call sequences for each prompt\footnote{Note that the use of cloud-based APIs does not pose privacy concerns in this context, as the training data can be generated within a controlled, dummy environment that excludes any sensitive information.}. We log the complete execution trajectories, which capture the initial prompt, the step-by-step LLM reasoning, the exact tool calls generated, and the resulting outputs from the tool. These high-quality trajectories serve as the ground truth for our fine-tuning process. While it may contain cases where \gptfivemini{} performs poorly, that such cases are rare, given the high quality of the frontier model. 
% \sean {but how we can be 100 percent sure the generated tool call sequence or trajectory by the frontier llm is correct?} \baris{yes, as we discussed this before, this is a weak point. need to somehow justify this.}. 
After generation, we partition the entire set of trajectories into an 80\% training set and a 20\% validation set. In addition, we reserve 50 randomly sampled queries as a separate test set (also referred to as \renabench{} which we evaluate in \sref{sec:evaluation}) that employs a balanced distribution to ensure coverage of all tools within each category. From each trajectory, we extract training instances for both the tool selection adapter and the relevant argument generation adapter. 

 % \keisuke{strengthen this argument}

\subsubsection{Separate Fine-Tuning for Tool Selection and Argument Generation}

% \Cref{??} shows the fine-tuning process for the generating the LoRA adapters in \sys{}. 
The training of the tool selection and argument generation adapters requires different inputs and applies different loss masking during fine-tuning, since they focus on different parts of a tool call. 

The tool selection adapter is trained on all the trajectories in the training set. The fine-tuning process involves predicting the name of the tool for the next step based on the context of all the prior steps and the user prompt for each trajectory. Since the aim of the tool selection adapter is only to choose the right tool (and not its arguments), we apply loss masking to compute loss only over the tool name tokens, not its arguments. 

% \sean{ this sentence implies we have one tool adapter per mcp server, is that the case? If yes, we need to probably discuss a bit earlier (e.g., we have only one tool adapter (or many? I cannot tell up to this point) but many argument adapters, and explain why we decided to choose such design) } 
The argument generation adapter of a particular tool is trained on all the appearances of the tool within the training trajectories. The fine-tuning process involves predicting the arguments for the tool, given the context of all prior steps in the trajectory and the tool name for the current step. The loss is computed over the tokens that contain the arguments for the tool call. 

% To ensure the adapters specialize correctly, we employ targeted loss masking during fine-tuning.

% When fine-tuning the classifier adapter, we mask the loss on all tokens except for those corresponding to the tool's name in the generated output. This forces the adapter to focus exclusively on learning the mapping from context to the correct tool choice, ignoring the complexity of argument generation.

% Conversely, when fine-tuning a tool-specific argument adapter, the loss masking is designed differently. \todo{@Zhan: Explain how the loss is masked to focus only on the argument structure and values, potentially masking the tool name itself.}\baris{is there some ablation result that measures the effect of loss masking?}

\subsection{Hierarchical Orchestration for Scalability}
As the number of available tools grows, a single classifier must choose from an increasingly large set of options, which increases context length and the complexity of the tool selection task. To address this, we introduce hierarchical orchestration to manage multiple toolsets (\ie, MCP servers). 
% \sean {maybe we can provide a definition of what is a tool what is a mcp server, what is the input output json schema earlier, it helps to understand the paper easier}\baris{yes even in the intro and at the begining of design, give an example first time use use the terms so it is clear.}.

This approach introduces a two-tiered tool selection process. At each step, we first use the \emph{base} local model without any fine-tuned adapters to perform a high-level routing task: selecting the most appropriate toolset (\eg, \texttt{filesystem}). We find that this high-level routing is more straightforward than fine-grained tool selection and can be effectively guided by a system prompt and structured decoding, without requiring fine-tuning. Intuitively, this is because the tools belonging to different toolsets are vastly different since they belong to different applications. Selecting an appropriate toolset is thus a simpler task than selecting a tool from \emph{within} a toolset. 
% \keisuke{could this be simply because we only have 3 toolsets? need to argue what if there are dozens of toolsets}
% \baris{provide some intuition as to why it doesn't require fine tuning}

Once a toolset is selected, we dynamically load the specialized tool selection adapter for that server. This tool selector then operates on a much smaller, more manageable set of tools (only the tools belonging to the selected toolset), significantly reducing the complexity of its decision. This hierarchical method effectively contains the context length and allows the system to scale to a larger number of tools without performance degradation. We evaluate hierarchical orchestration in \sref{sec:eval:ablation}.

% \baris{is there any eval of this. if so mention/fwd ref.}
% \baris{I think it is the first time I am realizing that there are in fact 3 layers of decision making and not two. I was under the impression that there is a classifier that selects the tool class and another that selects arguments, but this description says something else. I think this will confuse readers as we have led them to think the system architecture is otherwise. maybe a system diagram from the beginning would help, or somehow introducing this 3-tier architecture.}

\input{figures/fig-design-dualtune}

\subsection{Inference Framework}
% \sean{ if we could put this section earlier it helps to understand the system easier}

We integrate these techniques into the \sys{} inference framework, a complete agentic orchestration framework implemented in Rust. For security and stability, \sys{} executes all tools in containerized environments, providing fault isolation and resource management.

Our inference backend is powered by vLLM \cite{kwon2023efficient}, which is crucial for the practical feasibility of our approach. vLLM's ability to efficiently load, unload, and swap hundreds of LoRA adapters on the fly allows us to dynamically activate the necessary classifier and argument adapters at each step with minimal overhead. \Cref{fig:dualtune:design} shows the inference worflow in \sys{}:

\begin{enumerate}
    \item \textbf{Toolset Selection:} Given the current user prompt and conversation history, the base LLM selects the appropriate toolset (such as \texttt{filesystem} and \texttt{Notion}).
    \item \textbf{Tool Selection:} \sys{} loads the tool selection adapter for the selected MCP server. An inference call is made to this adapter, which returns the name of the desired tool.
    \item \textbf{Argument Generation:} \sys{} then dynamically swaps in the LoRA adapter specific to the chosen tool. A second inference call generates the precise arguments for the tool call.
    \item \textbf{Execution and Observation:} The generated tool call is executed in a secure container. The output is captured and appended to the conversation history.
    \item \textbf{Termination:} This loop continues until the tool selection adapter outputs a special token, ``summarize''. The system then generates a final summary for the user and terminates the execution.
\end{enumerate}
\section{Evaluation}
\label{sec:evaluation}

In this section, we present a comprehensive evaluation of different models and conduct ablation studies to analyze the contribution of individual components. We use \sys{} fine-tuning to create an LLM orchestrator called \sysmodel{}, which creates the LoRA adapters for \baselinemodel{}. 

\subsection{Experimental Setup}
\paragraph{Baselines}
We compare \sysmodel{} against local base models (\baselinemodel{}, Qwen-3-14B, Qwen-3-32B-Quant, Llama-3.1-3B), fine-tuned models (\toolace{}, \xlam{}) and frontier models (GPT-OSS-20b, GPT-OSS-120b, GLM-4.5, Kimi-K2-Instruct, DeepSeek-v3.1). 
% for these models are detailed in \Cref{??}.

\paragraph{Tools Evaluated}
Our evaluation encompasses 30 MCP tools spanning three distinct toolsets (detailed below), representing common real-world applications. We conduct 50 queries per toolset for each benchmark to ensure statistical reliability.

\begin{itemize}
    \item \texttt{Filesystem} toolset: This includes file and directory operations such as reading, writing, creating, and deleting files, as well as directory traversal and permission management tasks that form the backbone of system administration workflows.

    \item \texttt{monday.com} toolset: These tools encompass project management functionalities including task creation, status updates, team collaboration features, and workflow automation capabilities commonly used in enterprise project management environments.

    \item \texttt{Notion} toolset: This suite covers knowledge management operations including page creation, database queries, content organization, and collaborative editing features essential for modern documentation and information sharing workflows.
    
\end{itemize}

\paragraph{Benchmarks}
We assess \sys{} using two complementary benchmarks that evaluate different performance dimensions. The first benchmark is the test evaluation set of \sys{} (called \renabench{}), which includes 50 tasks for each toolset. The queries are designed to comprehensively test every tool in the toolset by generating roughly the same number of test cases for each tool, ensuring that all tools receive equal evaluation coverage. The second benchmark is MCP-Bench~\cite{wang2025mcp}, an open-source benchmark featuring predefined system prompts that generate complex, fuzzed queries requiring coordination across multiple diverse tools, though with non-uniform tool usage patterns when compared to \renabench, as explained in \sref{analysis:benchmark_metric}. We do not evaluate queries with writes, since we cannot judge objectively based on the content generated by the LLMs for writes. 

\paragraph{Evaluation Metrics}
Model performance is assessed using \judgemetric{}, a score ranging from 0 to 10 as detailed in \sref{analysis:benchmark_metric}. 
% We also report API costs for frontier models, which scale proportionally with input and output token consumption. 
Our evaluation employs GPT-5 as the judge model, which analyzes trajectory lists alongside ground-truth data to compute \judgemetric{} scores based on tool output quality and correctness.

\subsection{Overall Performance}

\input{tables/tbl-overall-perf}

\Cref{tbl:eval:perf-pivot} presents the comparative performance of all evaluated models across both benchmarks. We run all the models with the help of \sys{} and enable hierarchical orchestration for all the models. 
% Hence, all the results shown in this section reflect the improvements due to \decoupled{} used to create \sysmodel{}.

\paragraph{\sysmodel{} versus local models without reasoning}
\sysmodel{} achieves a higher accuracy compared to all the other local models (\baselinemodel{}, \xlam{} and \toolace{}) in both the benchmarks and across all the toolsets. 

\paragraph{\sysmodel{} versus reasoning models}
\sysmodel{} achieves a higher accuracy despite being a non-reasoning model compared to the local reasoning models (Qwen-3-14B, Qwen-3-32B-Quant) in 5 out of 6 cases.  

\paragraph{\sysmodel{} versus base model}
The \sys{} fine-tuning approach significantly improves the performance of \sysmodel{} compared to \baselinemodel{} on which it is post-trained, and achieves an accuracy up to 60\% higher than \basemodel{}. 

\subsection{Ablation Studies}
\label{sec:eval:ablation}
\input{figures/fig-eval-ablation}
To understand the individual contributions of \sys{}'s components, we conduct systematic ablation studies examining the effects of \decoupled{} and hierarchical orchestration.

We first begin with the base model (\baselinemodel{}) without hierarchical orchestration. To this we add hierarchical orchestration that reduces the context length of the model due to using a restricted set of tools. Then we test the performance of \baselinemodel{} with hierarchical orchestration and traditional fine-tuning, creating a single LoRA adapter using the training data generated in \sys{}. Finally, we evaluate \sysmodel{}, which is \baselinemodel{} with hierarchical orchestration and \decoupled{} instead of traditional fine-tuning. 

We perform this experiment on the \texttt{Filesystem} toolset for the queries generated using \mcpbench. The results for this ablation study are in \Cref{fig:eval:ablation}. We find that hierarchical orchestration improves the accuracy of \basemodel{} from 10.4\% to 16\%. Using traditional fine-tuning further improves the accuracy from 16\% to 39\%. However, when we substitute traditional fine-tuning with \decoupled{} for the same training dataset, we find that the accuracy jumps from 16\% to 61.5\%. 

% \vfill
\section{Limitations and Future Work}
\label{sec:limitations}
\sys{} relies on fine-tuning the base model to create a LoRA adapter to act as the tool selector for a toolset to select between available tools, and a separate LoRA adapter for each tool for argument generation. This has consequences in that adding new tools requires fine-tuning the common tool selector, as well as a separate LoRA adapter for the new tool. However, we believe that this is still a feasible approach because of two reasons. First, new tools are not added very frequently to existing tool sets. For example, there has only been one tool added to the \texttt{monday.com} toolset in the last five months, and no tools added to the \texttt{filesystem} toolset as well as the \texttt{Notion} toolset. This makes the fine-tuning process rare, and thus, practical. Second, we have a fully automated training pipeline as part of our contributions, which automatically generates the training dataset, performs the fine-tuning, and integrates the LoRA adapters as part of \sys{} -- the end-to-end process requires less than 10 hours in our experience, even while working with complex tool sets such as \texttt{monday.com} and \texttt{Notion}. In the future, we hope to explore techniques such as reinforcement learning to train the model while performing inference, thus automatically adapting to newer tools. 

\section{Related Works}

\subsection{LLMs on end-user devices}

Due to privacy concerns and the high costs of using advanced AI models through cloud APIs~\cite{hasan2025model, hou2025model, guo2025systematic}, researchers have explored ways to run AI models locally, including small foundation models~\cite{belcak2025small,apple_macbookpro_2024} and systems support for improving efficiency~\cite{frantar2022gptq, llama_cpp}. % These efforts have enabled AI inference on consumer-grade GPUs, making LLMs available to more people without relying on expensive cloud services.

\subsection{LLMs as agent orchestrators}

In recent years, there has been a significant interest in building autonomous agents based on LLMs that can execute complex tasks by calling an external tool. An increasing number of models, both in frontier models~\cite{openai2025introducingGPT5, anthropic2025introducingClaude4} and local models~\cite{zhang2024xlam,lin2024hammer,liu2024toolace}, support tool calling to function as agent orchestrators. The popularity of LLM agents has led to the creation of standard protocols, notably the Model Context Protocol (MCP)~\cite{anthropic_mcp_2024}
% and Agent2Agent Protocol \cite{google2025announcingA2A}. 
MCP allows applications to expose themselves as MCP Servers, containing a list of tools along with their descriptions. Using this protocol, LLMs can connect to external applications and issue tool calls to the connected toolsets, allowing them to interact with the applications autonomously.  

Additionally, there have been advancements in developing datasets, benchmarks, and leaderboards to assess the tool-calling capabilities of LLMs \cite{patil2025bfcl}. In the context of MCP, benchmarks like MCP-Bench are designed to assess the capability of LLMs to generate MCP tool-calls~\cite{wang2025mcp}, which involve realistic user prompts encompassing various applications and domains.

\subsection{Post Training for Tool Calling}

Some works propose post-training methods to improve tool calling capability. TinyAgent~\cite{erdogan2024tinyagent} performs fine-tuning with dataset generation that includes negative examples, and employs Retrieval Augmented Generation to select relevant tools corresponding to a user query, for reducing the context length of the model. Hammer~\cite{lin2024hammer} employs an augmented dataset that enhances models’ sensitivity to irrelevant functions and incorporates function masking techniques to minimize misleading. ToolACE~\cite{liu2024toolace} generates diverse tool-learning data to build a powerful local model. 

\section{Conclusion}
\label{sec:conclusion}
This paper proposes \emph{\decoupled{}}, a novel post-training approach that  enhances the tool calling capability of LLMs by disaggregating a tool-call into tool selection and argument generation sub-tasks, and separately fine-tunes these sub-tasks. This paper presents \sys{}, an LLM agent inference framework that uses LoRA adapters generated using \decoupled{} to perform efficient agent orchestation on consumer-grade GPUs. We plan to open-source our fine-tuning pipeline (\decoupled{}) and our inference framework (\sys{}). 

